% ===== S15 Directed Fissures =====
\section{Directed Fissures}
\label{sec:fissures}

\noindent
This paper does not conclude, and the refusal is principled rather than stylistic. A synthesis gathering the argument into a single closing proposition would settle what the argument has held open, converting the generative into the finished and enacting, in its own form, the appropriative closure the paper opposed at every register. In place of a conclusion the paper leaves directed fissures: places where what has been said runs out and points beyond itself, marked so that the pointing is legible. There are four, and each is a real incompleteness rather than a rhetorical modesty, an edge at which the argument reached the limit of what it could establish and can only indicate the direction of what it could not.

\subsection{The sedimentation of shared perception}
\label{sub:fiss-culture}

The first fissure is the widest, and the next paper must enter it. This paper showed how a shared everyday perception is generated between two people; it did not show what becomes of that perception when it accumulates. A shared perception generated and regenerated over the long iteration of a common life does not vanish at each turn. It sediments, into shared words, recurring gestures, the small rituals and understandings by which two people come to inhabit a perceptual world that is theirs and that a third could enter only by learning it. This sedimentation is the beginning of a culture, the smallest culture there is, generated in the intimacy of a shared everyday and carried beyond the moment of its generation into something inheritable. How the good repetition of shared perception sediments into a culture, how that culture is transmitted without being falsified, and how the smallest culture of two is continuous with the larger cultures of peoples and traditions, this paper has not shown. It is the subject of the paper that follows, on the emergence of culture in intimate relation, and this fissure is the seam along which the two are joined.

\subsection{The perception that cannot be shared}
\label{sub:fiss-untranslatable}

The second fissure is a limit inside the argument rather than a passage beyond it. The paper made the symbolic order the condition of shared perception, holding that a perception must be symbolised to be held in common. But there are perceptions that resist symbolisation, that remain in the private register of the imaginary or gather around a real so singular that no symbol carries them, and these the paper's account reaches only to mark as unreachable. What becomes of the aesthetic perception that cannot be shared, whether it is simply lost to the relational life the paper cares about or whether it plays some role the paper has not seen precisely in its resistance to sharing, is a question the argument raised and could not answer. The fissure is the edge of the shareable, and the paper can point to it without crossing it.

\subsection{The scale of the impersonal}
\label{sub:fiss-scale}

The third fissure concerns scale. The paper's unit was the relation of two, and its models were models of a shared life between persons who perceive each other perceiving. But perception is cultivated and starved at scales far larger than the dyad, by institutions, cities, and economies, and the paper's account of justice gestured at these scales without developing a practice for them. What generative relational aesthetic education would be at the scale of a school, a city, or a polity, how the models of a shared life between two could be extended to the impersonal relations of many without becoming the procedures the paper forbade, is undeveloped here. The fissure is the distance between the two who share an evening and the many who share a world, and the paper crossed it only in the direction of justice, not of practice.

\subsection{The perception of the last day}
\label{sub:fiss-finitude}

The fourth fissure is the one the paper is least able to close, and it opens where the everyday meets its end. The aesthetics of the everyday is an aesthetics of recurrence, of the light that returns each morning and the days that go on; its models assume the continuation of the shared life whose perception they cultivate. But shared lives end, in departure and in death, and the last ordinary day is not known to be the last while it is lived. What the generative perception of the everyday becomes under the knowledge of its ending, whether the finitude that the way of tea named in \zh{一期一会} deepens the everyday's generation or defeats it, and how a shared aesthetic life is to hold the fact that one of its two will perceive the other for the last time on a day that will have looked ordinary, the paper has felt at its edges and not thought through. This is the fissure that no next paper closes, and the paper leaves it open in the only way it can, by turning from theory to the thing itself.

\goldrule

\begin{center}
\begin{minipage}{0.82\linewidth}
\centering\itshape\color{rose}

{\cjk 晨光落案未成书，}\\[0.2em]
{\cjk 茶冷方知坐已虚。}\\[0.2em]
{\cjk 不必名之为美者，}\\[0.2em]
{\cjk 与君同看即吾庐。}\\[1.4em]

\upshape\normalcolor
Morning light falls on the desk; nothing yet is written.\\[0.15em]
The tea has gone cold, and only then do I know how long we have sat.\\[0.15em]
There is no need to give it the name of beauty.\\[0.15em]
To watch it with you is already my home.

\end{minipage}
\end{center}

\vspace{1em}

\noindent
The poem is the paper's last instantiation of its own claim, and it is offered in that spirit rather than as ornament. It symbolises a shared everyday perception without settling it, leaves the real value at its centre unnamed by declining the word \emph{beauty}, and returns the perception, through the reader, to the ordinary mornings it came from. It is a small generative symbolisation, made and shared, of exactly the kind the paper theorised, and it enacts, in its four lines, the whole of what the argument was for: that the light on an ordinary morning, watched with the one whose shared world it is, generates between two people a value that needs no name, accumulates into a home, and asks of theory only that it be seen, named where naming serves it, and defended where it is under threat. Here the theory ends and hands what it understood back to the ordinary life it was always about.
